\documentclass{article}
\usepackage[T1]{fontenc}
\usepackage{lmodern}
\usepackage{graphicx}
\usepackage{amsmath,amssymb}
\usepackage[a4paper,margin=2cm]{geometry}
\usepackage[absolute,overlay]{textpos}
\setlength{\TPHorizModule}{1mm}
\setlength{\TPVertModule}{1mm}
\usepackage[indonesian]{babel}
\usepackage{hyperref}
\setlength{\emergencystretch}{2em}
% unit: Halaman 10

\begin{document}

% === Halaman 10 (python/native) ===

• Kemudian pesan 32-bit ini digeser ke 

kiri sejauh 11 bit secara sirkuler. 

• Hasilnya kemudian di-XOR-kan dengan 

Li – 1  untuk kemudian memberikan 

bagian cipherteks kanan yang baru, Ri. 

Proses ini diulang sebanyak 32 kali. 

• Dekripsi merupakan kebalikan dari 

enkripsi. Dari “bawah” ke “atas” 

mengikuti jaringan Feistel. 

Li - 1

Ri - 1

Substitusi

Kotak-S

Pergeseran sirkuler

ke kiri 11 bit



Ri

Li

Ki

Keterangan:

adalah operator penjumlahan dalam modulo 232

f

10

\end{document}
